\chapter{前言}\label{preface}

\paragraph{关于本书}
正如书名所示，这是一本关于形式逻辑的教科书。形式逻辑的研究对象是一种特定的语言，它和任何语言一样，可以用来表达事态。它是一种形式语言，即其表达式（例如语句）是形式化定义的。这使得它成为一种非常有用的语言，能够极其精确地描述其语句所表达的事态。特别是在形式逻辑中，不可能出现歧义。对这些语言的研究，核心在于语句之间的蕴涵关系，即哪些语句可以从哪些其他语句得出。蕴涵之所以至关重要，是因为通过更好地理解它，我们就能判断，在某些事态成立时，另一些事态必然也成立。但蕴涵并非唯一重要的概念。我们还将考虑可满足关系，即语句间互不矛盾的关系。这些概念可以从语义角度来定义，即基于对语言的解释来精确定义蕴涵；也可以从证明论角度来定义，即使用形式化的演绎系统。

形式逻辑是哲学的一个核心分支学科，其中，前提与结论之间的逻辑关系至关重要。哲学家探究定义和假设的后果，并基于这些后果来评估这些定义和假设本身。形式逻辑在数学和计算机科学中也很重要。在数学中，形式语言被用来描述的不是“日常”事态，而是数学事态。数学家同样关心定义和假设的后果，并且对于他们而言，使用完全精确和严谨的方法来确立这些后果（即所谓的“定理”）也同样重要。形式逻辑提供了这样的方法。在计算机科学中，形式逻辑被应用于描述计算系统（例如电路、程序、数据库等）的状态和行为。形式逻辑的方法同样可以用来确立这类描述的后果，例如电路是否无误、程序是否实现了预期功能、数据库是否保持一致，或者关于其中数据的某个断言是否成立。

本书分为九个部分。\Cref{ch.intro}以非形式化的方式介绍逻辑这一主题及其相关概念，尚未引入形式语言。\Crefrange{ch.TFL}{ch.NDTFL} 讨论真值函项逻辑的语言。在此类语言中，语句通过若干联结词（“或”、“且”、“非”、“如果……那么”）从基本语句构成，这些联结词将语句组合成更复杂的语句。我们用两种方式讨论诸如蕴涵这样的逻辑概念：语义上，使用真值表方法（在\cref{ch.TruthTables}中）；证明论上，使用一个形式推导系统（在\cref{ch.NDTFL}中）。\Crefrange{ch.FOL}{ch.NDFOL} 处理一种更复杂的语言，即一阶逻辑。它除了包含真值函项逻辑的联结词外，还包括名称、谓词、等词以及所谓的量词。语言的这些额外要素使其表达力远超真值函项逻辑的语言，我们将花费相当篇幅探究在其中究竟能表达什么。同样，一阶逻辑语言的逻辑概念可以从语义角度（使用解释）和证明论角度（使用\cref{ch.NDTFL}中介绍的形式推导系统的一个更复杂的变体）来定义。\Cref{ch.ML}讨论对真值函项逻辑的扩展，即引入表达可能性与必然性的非真值函项算子：模态逻辑。\Cref{ch.normalform}涵盖两个高级主题：合取范式与析取范式以及真值函项联结词的功能完备性，以及真值函项逻辑自然演绎的可靠性。

在附录中，您将找到关于本书所讨论语言的替代记法、替代推导系统的探讨，以及一份列出了大部分重要规则和定义的速查参考。核心术语列于书末的术语表中。

\paragraph{版权说明} 本书基于 P.~D. Magnus 撰写、Tim Button 修订和扩充的版本。它还包含了 J.~Robert Loftis 提供的一些材料（主要是练习题）。\cref{ch.ML} 的材料基于 Robert Trueman 的笔记（但已改写为使用 Fitch 原始的模态逻辑自然演绎规则），而 \cref{c:NormalForms,c:FunctionalCompleteness,ch:Soundness} 的材料则基于 Tim Button 的开放教材《元理论》中的两章。Aaron Thomas-Bolduc 和 Richard Zach 将这些文本的内容整合为当前版本，更改了一些术语和例子，重写了一些章节，并加入了他们自己的材料。特别是，Richard Zach 重写了 \cref{s:Arguments,s:Valid}，并增补了 \cref{s:AmbiguityTFL,s:stratTFL,s:ambiguityFOL,ch:PropRelations,ch:equivalences}。截至 2019 年秋季版，关于一阶逻辑的部分采用了更常见于高级教材的语法（例如那些基于 \href{https://openlogicproject.org/}{开放逻辑项目}的教材），即谓词符号的变元置于括号内（例如用 “$R(a,b)$” 代替 “$Rab$”）。这一版本的 \forallx{} 也采用了一阶逻辑的标准语法定义（允许空量化），并使用 Gentzen 的原始自然演绎引入与消去规则（即双重否定消去规则和排中律规则是派生规则，间接证明是唯一的经典规则）。最终文本根据知识共享署名 4.0 许可协议发布。存在 \href{https://github.com/OpenLogicProject/OpenLogic/wiki/Other-Logic-Textbooks}{其他几个} \forallx 的“改编”版本，包括本版本的翻译。

\paragraph{致教师} 本书内容适合作为一学期的形式逻辑导论课程教材。我曾在 12 周内讲授 \crefrange{ch.intro}{ch.NDFOL} 以及 \cref{c:NormalForms,c:FunctionalCompleteness,ch:equivalences}，但会略去部分真值表与派生推理规则。

本书最新版本的 PDF 可在 \href{https://forallx.openlogicproject.org}{forallx.openlogicproject.org} 获取，但更新频繁。知识共享署名许可授予您下载和自行分发本书的权利。为确保您授课期间的所有学生使用的是同一版本，建议您将决定使用的 PDF 上传至您的学习管理系统，而不是仅仅提供链接。本书也提供 \href{https://forallx.openlogicproject.org/html/}{HTML 版本}。该版本在\href{https://forallx.openlogicproject.org/html/A4.html}{无障碍访问方面经过了特别考量}，应更适合依赖屏幕阅读器的用户，例如低视力或全盲的学生。您可以下载包含 HTML 和 PDF 版本的 SCORM 资源包，以便上传至您的学习管理系统。您也可以自行联系书店打印 PDF，但有些书店可能可以采购并库存亚马逊上有售的平装本。

请注意，书中许多练习的答案同样可在上述网站获取（面向所有人，包括您的学生）。这些答案（目前）尚未包含在 HTML 版本或 SCORM 资源包中。

本书的语法与证明系统（模态逻辑部分除外）由 Graham Leach-Krouse 的免费在线逻辑教学软件 \textit{Carnap}（\href{https://carnap.io}{carnap.io}）提供支持。这使得符号化、真值表构造和自然演绎证明等练习的提交与自动批改成为可能。Carnap.io 上的教师可以找到\href{https://carnap.io/shared/rzach@ucalgary.ca/forall%20x:%20Calgary.md}{可供改编或直接布置的附加练习示例}。
